\subsection{nTH ROOTS}

Let \(n\) be any integer \(> 0\) . From the relation \(0 < x < y\) we deduce, by induction on \(n\) , that \(0 < x^n < y^n\) . In other words, the function \(x \to x^n\) is strictly increasing for \(x \geqslant 0\) ; it is clearly continuous at every point and therefore (§ 2, no. 6, Theorem 5) it is a homeomorphism of \(\mathbf{R}_{+}\) onto an interval \(\mathbf{I}\) . On the other hand, since \(x \geqslant 1\) implies \(x^{n-1} \geqslant 1\) and therefore \(x^n \geqslant x\) , it follows that \(\mathbf{I}\) is not bounded and hence \(\mathbf{I} = \mathbf{R}_{+}\) . The value, for \(x \geqslant 0\) , of the inverse of the mapping \(x \to x^n\) is denoted by \(x^{1/n}\) or \(\sqrt[n]{x}\) and is called \(x\) to the power \(1\) or the nth root of \(x\) (for \(n = 2\) , 3 we say square root, cube root; for \(n = 2\) we write \(\sqrt{x}\) in place of \(\sqrt[2]{x}\) ). The positive number \(x^{1/n}\) is thus defined as the unique positive solution of the equation

\[
y ^ {n} = x \quad (x \geqslant 0).\tag{2}
\]

In particular we see that there is a real number x such that \(x^{2}=2\) , whereas no rational number has this property; thus we recover the fact that the rational line Q is not a complete space.

The mapping \(x \to x^{1/n}\) of \(\mathbf{R}_+\) onto itself is strictly increasing and continuous. By (2) we have \(0^{1/n}=0, 1^{1/n}=1\) , also

\[
(x y) ^ {1 / n} = x ^ {1 / n} y ^ {1 / n};\tag{3}
\]

hence \(x \to x^{1/n}\) is an automorphism of the topological group \(\mathbf{R}_{+}^{*}\) .

In Chapter V, § 4, no. 1, we shall generalize this result by finding all the automorphisms of the multiplicative group \(\mathbf{R}_{+}^{*}\) .

